% !TEX program = lualatex
\documentclass{../armymemo}
% used to show page layouts
\usepackage{layouts}
% In order to show baselines:
% https://tex.stackexchange.com/q/159010
\usepackage{atbegshi,picture,xcolor,hyperref}
\AtBeginShipout{%
  \AtBeginShipoutUpperLeft{%
    \color{red}%
    \put(\dimexpr 1in+\oddsidemargin,
    -\dimexpr \topmargin+\headheight+\headsep+\topskip)%
    {%
      \vtop to\dimexpr\vsize+\headheight+\headsep+\topskip+\baselineskip{
        \hrule
        \leaders\vbox to\baselineskip{\hrule width\hsize\vfill}\vfill
      }%
    }%
  }%
}

\logo{imperialseal-536x536}
\department{Galactic Fleet}
\address{Death Star}
\address{Alderaan System}
\address{Core Worlds}
\officesymbol{DS-1-OBS}
\signaturedate{10 April 2019}
\documentmark{GALACTIC CONFIDENTIAL//DEATHSTAR}

% If the subject spills into a 2nd line, it breaks \pagedesign from layout, but otherwise seems to work
\subject{On the use and testing of \LaTeX\ for document production}% and system documentation for the DS-1 Orbital Battle Station}

% \mfr

\multimemothru{U.S. Army Command and General Staff College (ATZL), 100 Stimson Avenue, Ft Leavenworth, KS 66027-1352}
\multimemothru{CHIEF OF STAFF, 501\st\ Legion, Galactic Army, Grand Army of the Republic, Galactic Empire}
\memoline{FOR General of the Grand Army of the Republic, Dark Lord of the Sith, Darth Vader}

% \memoline{MEMORANDUM THRU CHIEF OF STAFF, 501\st\ Legion, Galactic Army, Grand Army of the Republic, Galactic Empire}
% \memoline{FOR U.S. Army Command and General Staff College (ATZL), 100 Stimson Avenue, Ft Leavenworth, KS 66027-1352}

% \memoline{MEMORANDUM FOR U.S. Army Command and General Staff College (ATZL), 100 Stimson Avenue, Ft Leavenworth, KS 66027-1352}

\authority{FOR THE COMMANDER}
\author{Boba Fett}
\rank{Contractor}
\branch{Independent}
\title{Procurement of Prized Heads}

\suspensedate{1 JAN 2019}

\addencl{\LaTeX HOWTO}
\addencl{StackExchange}
\addencl{armymemo.cls}
\adddistro{G.M. Tarkin}
\adddistro{Adm. Ozzel}
\adddistro{Adm. Motti}
\adddistro{V.Adm. Piett}
\addcf{Galactic J3}
\addcf{Fleet J3}

\begin{document}
\begin{enumerate}
	% \item \currentlist
	%   \begin{figure}
	%     \listdiagram
	%   \end{figure}
	\item This document uses code from \url{https://tex.stackexchange.com/q/159010}
	      to display the baselines in order to confirm spacing between paragraph, list
	      items, and other elements of the memo.
	\item This \verb!example-grid.tex! document is for testing features of the class
	      and validating layout.
	\item See paragraph 2-2 (of AR 25-50) for when to use a memorandum.
	\item Single space the text and double space between paragraphs and
	      subparagraphs. Insert two blank spaces after ending punctuation (period and
	      question mark). Insert two blank spaces after a colon. When numbering
	      subparagraphs, insert two blank spaces after parentheses.
	\item When a memorandum has more than one paragraph, number the paragraphs
	      consecutively. When paragraphs are subdivided, designate first subdivisions
	      using lowercase letters of the alphabet and indent 1/4 inch as shown below.
	      \begin{enumerate}
		      \item When a paragraph is subdivided, it must have at least two subparagraphs.
		      \item If there is a subparagraph ``a,'' there must be a subparagraph ``b.''
		            \begin{enumerate}
			            \item Designate second subdivisions by numbers in parentheses; for example
			                  (1), (2), and (3) and indent by 1/2 inch as shown.
			            \item Do not subdivide beyond the third subdivision.
			                  \begin{enumerate}
				                  \item Do not indent any further than the second subdivision.
				                  \item Use (a), (b), (c), and so forth at this level.
			                  \end{enumerate}
		            \end{enumerate}
	      \end{enumerate}
	\item Moby Dick, a classic American Novel, the text of which is public domain,
	      follows in several paragraphs to make a longer example:
	      \begin{enumerate}
		      \item Call me Ishmael.\footnote{This is a footnote} Some years ago- never mind
		            how long precisely- having little or no money in my purse\footnote{wallet},
		            and nothing particular to interest me on shore, I thought I would sail about
		            a little and see the watery part of the world.\footnote{aka the ocean} It is
		            a way I have of driving off the spleen and regulating the
		            circulation. Whenever I find myself growing grim about the mouth; whenever
		            it is a damp, drizzly November in my soul; whenever I find myself
		            involuntarily pausing before coffin warehouses, and bringing up the rear of
		            every funeral I meet; and especially whenever my hypos get such an upper
		            hand of me, that it requires a strong moral principle to prevent me from
		            deliberately stepping into the street, and methodically knocking people's
		            hats off- then, I account it high time to get to sea as soon as I can.  This
		            is my substitute for pistol and ball. With a philosophical flourish Cato
		            throws himself upon his sword; I quietly take to the ship. There is nothing
		            surprising in this. If they but knew it, almost all men in their degree,
		            some time or other, cherish very nearly the same feelings towards the ocean
		            with me.

		      \item There now is your insular city of the Manhattoes, belted round by wharves as
		            Indian isles by coral reefs- commerce surrounds it with her surf. Right and
		            left, the streets take you waterward. Its extreme downtown is the battery, where
		            that noble mole is washed by waves, and cooled by breezes, which a few hours
		            previous were out of sight of land. Look at the crowds of water-gazers there.

	      \end{enumerate}

	\item POC for this memo is John Doe, who can be reached at \texttt{john@dev.null}
	      or (555) 555-1234.
\end{enumerate}

\begin{figure}
	\centering
	\currentpage
	\drawmarginparsfalse
	\drawdimensionstrue
	\pagedesign
	\caption{Document layout}
	\label{fig:1}
\end{figure}

\end{document}